\pagestyle{fancy}
\fancyhead{}

\fancyheadoffset{0cm}
\renewcommand{\headrulewidth}{1pt} 
\renewcommand{\footrulewidth}{0pt}
\fancyhead[L]{{\thesection}}
\fancyhead[L]{\leftmark}

% \fancyhead[LO,LE]{\nouppercase\firstleftmark}
%\fancyhead[RO,RE]{\thesubsection\enspace\subsectiontitle}

\section{Conclusions}

This study proposes an approach to generalize SOZ localization to different
patients using a unique set of parameters.

The proposed pipeline obtained good overall performances on training and
validation patients, however all evaluation metrics appear highly variable,
highlighting the challenge of finding universal parameters to fit such a
diverse dataset.
Despite obtaining promising results during cross-validation, the evaluation of
unseen patients resulted in a substantial drop in performances.
Major limitations were also observed for the Epileptogenicity
Index proposed by Bartolomei et al.:
although being a widely established algorithm for SOZ localization on
single-patient settings, such model fails to correctly classify a meaningful 
number of SOZ regions when the same parameters are applied to multiple patients.

This behavior suggests that, although the proposed pipeline still leaves room 
for improvement, the primary challenge of the task lies in the intrinsic nature
of the dataset.
sEEG data inherently presents inter-patient heterogeneity, both because of
subjective variability in brain activity and because of differences in the 
positioning of the electrodes during data acquisition.
The number of available patients for this study was also moderately low,
limiting generalization capabilities on unseen patients.
A larger set of patients would be preferable for this reaserch, however
potential strategies have been proposed to account for the problem of intrinsic
data variability.

Despite these challenges, the proposed framework provides an initial approach
to the SOZ localization of unseen patients in focal epilepsy.
While still not ready for clinical use, directions for future improvement have
been identified for a more robust and patient-adaptive approach.
